\documentclass[11pt,leqno,final]{amsart}

\usepackage{packages}

\title{MAC0325/5781 Combinatorial Optimization: Assignment 1}

\droptag{comments}
\usetag{hint}
\droptag{attribution}

\begin{document}

\maketitle

\begin{tabbing}
  \textbf{Instructions:} \= \kill
  \textbf{Submit:} \> \cref{ex:sptree-ip,ex:maxcut-ip,ex:odd-circuit-blocker}
  \\
  \textbf{Deadline:} \> 5/May/2024 at 11:59pm
  \\
  \textbf{Instructions:} \> Follow \textbf{CAREFULLY} the submission
  instructions listed in Moodle.
\end{tabbing}

\begin{exercise}
  Let \(G \coloneqq (V,E,\varphi)\) be a graph.
  Let \(F \subseteq E\).
  Prove that the following are equivalent:
  \begin{enumerate}[(i)]
  \item \(F\) is a spanning tree of~\(G\);
  \item \((V,F,\varphi\restrict{F})\) is connected and \(\card{F} = n-1\);
  \item \((V,F,\varphi\restrict{F})\) is connected and has no circuits
  \item \((V,F,\varphi\restrict{F})\) has no circuits and \(\card{F} = n-1\);
  \item \((V,F,\varphi\restrict{F})\) is a minimal spanning subgraph
    of~\(G\) that is connected.
  \end{enumerate}
  Also prove that, if \(G\) is connected, the following item may be
  added to the list:
  \begin{enumerate}[resume*]
  \item \((V,F,\varphi\restrict{F})\) is a maximal subgraph of~\(G\)
    that has no circuits.
  \end{enumerate}
\end{exercise}

\section{Integer Programming Formulations}

\begin{exercise}
  \label{ex:atsp}
  Let \(D \coloneqq (V,A,\varphi)\) be a digraph, and let
  \(c \colon A \to \Reals_+\) be a function that assigns a nonnegative
  ``cost'' to each arc.
  A circuit~\(C\) in~\(D\) is \textbf{Hamiltonian} if~\(V(C) = V\).
  Formulate the problem of finding a minimum-cost Hamiltonian circuit
  of~\(D\) as an integer program, and argue that the formulation is
  correct.
  The cost of a Hamiltonian circuit~\(C\)
  is~\(\iprodt{c}{\incidv{A(C)}}\).
\end{exercise}

\begin{exercise}
  Let \(D \coloneqq (V,A,\varphi)\) be a digraph, and let
  \(c \colon A \to \Reals_+\) be a function that assigns a nonnegative
  ``cost'' to each arc.
  Let \(r \in V\).
  Prove correctness that the following IP formulation
  is correct for \cref{ex:atsp}:
  \begin{alignat*}{3}
    \text{Minimize} \quad
    & \iprodt{c}{x} &
    \\[0pt]
    \text{subject to}\quad
    & x \in \set{0,1}^A, &&
    \\[0pt]
    & x\paren[\big]{\deltaout(i)} = 1, &&
    \qquad \text{for each }i \in V,
    \\[0pt]
    & x\paren[\big]{\deltain(i)} = 1, &&
    \qquad \text{for each }i \in V,
    \\[0pt]
    & u \in \Integers^{V \drop \set{r}}, &&
    \\[0pt]
    & u_i - u_j + n x_{ij} \leq n-1, &&
    \qquad \text{for each }ij \in A\text{ disjoint from }\set{r}.
  \end{alignat*}
\end{exercise}

\begin{exercise}
  \label{ex:sptree-ip}
  Let \(G \coloneqq (V,E)\) be a graph, and let
  \(c \colon E \to \Reals_+\) be a function that assigns a nonnegative
  ``cost'' to each edge.
  Formulate the problem of finding a minimum-cost spanning tree
  of~\(G\) as an integer program, and argue that the formulation is
  correct.

  \smallskip
  \noindent\small
  \textit{Page Limit:} 1 page.
\end{exercise}

\begin{exercise}
  \label{ex:sptree-neg}
  Let \(G \coloneqq (V,E)\) be a graph, and let
  \(c \colon E \to \Reals\) be a function that assigns a (possibly negative)
  ``cost'' to each edge.
  Formulate the problem of finding a minimum-cost spanning tree
  of~\(G\) as an integer program, and argue that the formulation is
  correct.
\end{exercise}

\begin{exercise}
  \label{ex:maxcut-ip}
  Let \(G \coloneqq (V,E)\) be a graph.
  Recall that a \textbf{cut} of~\(G\) is a subset of~\(E\) of the form
  \begin{equation*}
    \delta_G(S)
    \coloneqq
    \setst{
      e \in E
    }{
      e \cap S \neq \emptyset \neq
      e \cap \overline{S}
    }
  \end{equation*}
  for some \(S \subseteq V\), where \(\overline{S}\) denotes
  \(V \drop S\).
  Formulate the problem of finding a cut of maximum size in a graph as
  an integer program, and argue that the formulation is correct.

  \smallskip
  \noindent\small
  \textit{Page Limit:} 1 page.

  \tagged{hint}{%
    \smallskip

    \noindent\small
    \emph{Hint.}
    It might help to add extra variables for each vertex and edge.
  }
\end{exercise}

\begin{exercise}
  Let \(G \coloneqq (V,E)\) be a graph.
  Formulate the problem of finding a minimum(-size) coloring of~\(G\)
  as an integer program, and argue that the formulation is correct.
\end{exercise}

\begin{exercise}
  Let \(G \coloneqq (V,E)\) be a graph.
  A dominating set \(D \subseteq V\) of~\(G\) is \textbf{connected} if
  the induced subgraph \(G[D]\) is connected.
  Formulate the problem of finding a minimum(-size) connected
  dominating set of~\(G\) as an integer program, and argue that the
  formulation is correct.
\end{exercise}

\section{Clutters and Blockers}

Recall the following basic notions of hypergraphs.
A \textbf{hypergraph} is an ordered pair
\(\Hcal \coloneqq (V,\Ecal)\), where \(V\) is a finite set whose
elements are called the \textbf{vertices} of~\(\Hcal\) (and denoted by
\(V(\Hcal)\) or \(V_{\Hcal}\)), and \(\Ecal \subseteq \Powerset{V}\).
The elements of \(\Ecal\) are called \textbf{hyperedges} (or just
\textbf{edges}) of~\(\Hcal\), and \(\Ecal\) is also denoted by
\(\Ecal(\Hcal)\) or \(\Ecal_{\Hcal}\).
If \(\Hcal_1\) and \(\Hcal_2\) are hypergraphs, write
\(\Hcal_1 \leq \Hcal_2\) to mean that
\(V(\Hcal_1) \subseteq V(\Hcal_2)\) and
\(\Ecal(\Hcal_1) \subseteq \Ecal(\Hcal_2)\).
For each hypergraph \(\Hcal \coloneqq (V,\Ecal)\), define the
hypergraphs
\begin{align*}
  \Hcal^{\min}
  & \coloneqq
  \paren[\Big]{
  V,
  \setst[\big]{F \in \Ecal}{\text{there is no \(E \in \Ecal\) such that
    \(E \subsetneq F\)}}
  },
  \\
  \Hcal^{\uparrow}
  & \coloneqq
  \paren[\Big]{
  V,
  \setst[\big]{F \subseteq V}{\text{there is an \(E \in \Ecal\) such that
    \(E \subseteq F\)}}
  }.
\end{align*}

Hypergraphs are an important source of IP formulations of
combinatorial optimization problems.
Let \(\Hcal \coloneqq (V,\Ecal)\) be a hypergraph.
A \textbf{vertex cover} of~\(\Hcal\) is a subset \(K \subseteq V\)
such that, for each hyperedge \(E \in \Ecal\), one has
\(K \cap E \neq \emptyset\).
For each \(w \in \Reals_+^V\), finding a minimum-weight vertex cover
of~\(\Hcal\) can be formulated as the following IP:
\begin{alignat*}{3}
  \text{Minimize} \quad
  & \iprodt{w}{x} &
  \\[0pt]
  \text{subject to}\quad
  & x \in \set{0,1}^V, &
  \\[0pt]
  & x(F) \geq 1, & \text{ for each }F \in \Ecal.
\end{alignat*}

A hypergraph \(\Hcal \coloneqq (V,\Ecal)\) is a \textbf{clutter} if
\(\Ecal\) is an antichain with respect to inclusion, i.e., if
\(\Hcal^{\min} = \Hcal\).
The \textbf{blocker} of a hypergraph \(\Hcal \coloneqq (V,\Ecal)\) is
the clutter
\begin{equation*}
  b(\Hcal)
  \coloneqq
  \paren[\big]{
    V,
    \setst{K \subseteq V}{K \text{ is a vertex cover of~\(\Hcal\)}}
  }^{\min}.
\end{equation*}
Thus, the hyperedges of \(\paren[\big]{b(\Hcal)}^{\uparrow}\) are
precisely the vertex covers of~\(\Hcal\).
A~very important (combinatorial) duality relation, due to Lawler and,
independently, Edmonds and Fulkerson, is that, if \(\Hcal\) is a
clutter, then \(b(b(\Hcal)) = \Hcal\).
A clutter and its blocker are said to be a \textbf{blocking pair} of
clutters.

\begin{exercise}
  Let \(\Hcal\) and \(\Bcal\) be hypergraphs on the same vertex set,
  such that
  \begin{equation*}
    \Bcal \leq \paren[\big]{b(\Hcal)}^{\uparrow}
    \qquad\text{and}\qquad
    \paren[\big]{b(\Bcal)}^{\uparrow} \leq \Hcal^{\uparrow}.
  \end{equation*}
  Prove that \(b(\Hcal) = \Bcal^{\min}\), so \(\Hcal^{\min}\) and
  \(\Bcal^{\min}\) form a blocking pair of clutters.
\end{exercise}

\begin{exercise}
  \label{ex:tree-cut-blocker}
  Let \(G = (V,E)\) be a graph.
  Define the hypergraphs
  \begin{align*}
    \Hcal
    & \coloneqq
    \paren[\big]{
      E,
      \setst{T \subseteq E}{T \text{ is a spanning tree of~\(G\)}}
    },
    \\
    \Bcal
    & \coloneqq
    \paren[\big]{
      E,
      \setst{\delta(S)}{S \in \Powerset{V},\, \emptyset \neq S \neq V}
    }.
  \end{align*}
  Prove that \(\Hcal^{\min}\) and \(\Bcal^{\min}\) form a blocking
  pair of clutters.
\end{exercise}

\begin{exercise}
  \label{ex:rs-path-blocker}
  Let \(G = (V,E)\) be a graph, and let \(r,s \in V\) be distinct.
  Define the hypergraphs
  \begin{align*}
    \Hcal
    & \coloneqq
    \paren[\big]{
      E,
      \setst{E(P)}{P \text{ is an \(rs\)-path in~\(G\)}}
    },
    \\
    \Bcal
    & \coloneqq
    \paren[\big]{
      E,
      \setst{\delta(S)}{S \in \Powerset{V \setminus \set{t}},\, s \in S}
    }.
  \end{align*}
  Prove that \(\Hcal^{\min}\) and \(\Bcal^{\min}\) form a blocking
  pair of clutters.
\end{exercise}

\begin{exercise}
  \label{ex:odd-circuit-blocker}
  Let \(G = (V,E)\) be a graph.
  You may assume that \(G\) is not bipartite, though what follows
  still holds for bipartite graphs.
  Define the hypergraphs
  \begin{align*}
    \Hcal
    & \coloneqq
    \paren[\big]{
      E,
      \setst{E(C)}{C \text{ is an odd circuit of~\(G\)}}
    },
    \\
    \Bcal
    & \coloneqq
    \paren[\big]{
      E,
      \setst{E \setminus \delta(S)}{S \in \Powerset{V},\, \emptyset \neq S \neq V}
    }.
  \end{align*}
  Prove that \(\Hcal^{\min}\) and \(\Bcal^{\min}\) form a blocking
  pair of clutters.

  \smallskip
  \noindent\small
  \textit{Page Limit:} 1 page.
\end{exercise}

\begin{exercise}
  \label{ex:Tjoin-blocker}
  Let \(G = (V,E)\) be a graph, and let \(T \subseteq V\) such that
  \(\card{T}\) is even.
  A subset \(J \subseteq E\) is a \textbf{\(T\)-join} if every vertex
  of \((V,J)\) has even degree.
  Define the hypergraphs
  \begin{align*}
    \Hcal
    & \coloneqq
    \paren[\big]{
      E,
      \setst{J \subseteq E}{J \text{ is a \(T\)-join of~\(G\)}}
    },
    \\
    \Bcal
    & \coloneqq
    \paren[\big]{
      E,
      \setst{\delta(S)}{S \in \Powerset{V},\, \card{S \cap T} \text{
      is odd}}
    }.
  \end{align*}
  Prove that \(\Hcal^{\min}\) and \(\Bcal^{\min}\) form a blocking
  pair of clutters.
\end{exercise}

\end{document}

%%% Local Variables:
%%% mode: latex
%%% TeX-master: t
%%% coding: utf-8
%%% sentence-end-double-space: t
%%% End:
